% ==========================================================
% 4. Minimundo do Projeto (Contexto atual)
% ==========================================================
\chapter{Minimundo do Projeto (Contexto atual)}
\label{cap:minimundo}

A Commit Júnior é uma Empresa Júnior (EJ) sem fins lucrativos, gerida exclusivamente por discentes de cursos de graduação em Engenharia do Centro Federal de Educação Tecnológica de Minas Gerais (CEFET-MG). A organização atua no mercado fornecendo serviços de consultoria e soluções computacionais, estruturando seu quadro operacional e administrativo em diretorias e gerências especializadas, tais como Diretoria Executiva, Gestão de Pessoas, Jurídico, Comercial e Projetos. Para a realização das atividades diárias, os membros utilizam notebooks, computadores pessoais e smartphones, mantendo a comunicação por meio de redes Wi-Fi acadêmicas nos campi da instituição de ensino, redes residenciais em regime remoto e pacotes de dados móveis com conectividade estável. O corpo de membros é composto por graduandos de 18 a 26 anos, todos plenamente alfabetizados, comunicativos e dotados de alto nível de proficiência técnica no uso de tecnologias digitais e serviços em nuvem.

No fluxo operacional vigente, o gerenciamento de pessoas, a formalização contratual e o controle das rotinas internas dependem de processos descentralizados e manuais, divididos entre diferentes ferramentas e responsáveis:

No âmbito cadastral e contratual, a captação inicial dos dados pessoais e civis dos membros é feita por formulários do Google Forms. Essas respostas são exportadas para planilhas do Google Sheets gerenciadas pelo setor Jurídico. A partir daí, o Jurídico precisa consultar individualmente as linhas da planilha, conferir os dados e gerar manualmente os termos de voluntariado a partir de um modelo padrão, os quais são posteriormente encaminhados para assinatura eletrônica por meio da plataforma Autentique. O controle do retorno dessas assinaturas é registrado manualmente na planilha de membros por meio de códigos de verificação de autenticidade dos termos. O mesmo procedimento manual e moroso se repete quando ocorre a saída de um integrante: o Jurídico redige um termo de desligamento nos mesmos moldes, colhe as assinaturas via Autentique e atualiza o status na planilha. De forma análoga, a emissão e o acompanhamento de certificados de participação na empresa júnior passam por essa conferência manual de registros e preenchimento de modelos para cada membro.

No gerenciamento de extensão acadêmica, a contabilização das horas é mantida em uma planilha específica do Google Sheets. Ao término do ciclo de permanência do membro na empresa, o assessor da presidência realiza o cálculo manual do tempo de casa e da quantidade total de horas de extensão acumuladas em cada cargo desempenhado. Esse cálculo manual toma muito tempo e está sujeito a erros contínuos de contagem e digitação. Concluída a soma, as informações do aluno e o total de horas são enviados de maneira informal pelo WhatsApp para o professor orientador da Commit Júnior. O professor orientador, por sua vez, realiza a inserção desses dados no SIGAA, no qual é lançado o processo de extensão; a partir disso, o membro precisa elaborar e submeter um relatório individual de atividades na plataforma acadêmica para que, após a devida aprovação, receba o certificado de extensão oficial do CEFET-MG.

Paralelamente, a consulta de contatos internos, squad e cargos dos integrantes em exercício é mantida em quadros colaborativos no Notion, ferramenta na qual a liderança também preenche um painel de lançamentos de pontos para controle disciplinar (faltas, advertências e condutas inadequadas). A verificação dos prazos de validade dessas pontuações e o cômputo dos agravantes por reincidência exigem conferência estritamente visual dos gestores, pois a ferramenta não possui automação para caducidade nem para multiplicação de penalidades. Por fim, as Reuniões Gerais (RG) e os alinhamentos bilaterais (1-on-1) são organizados de modo descentralizado por aplicativos de mensagem instantânea e calendários pessoais, resultando frequentemente em sobreposições de agenda e choques com compromissos corporativos obrigatórios.

Essa dinâmica fragmentada gera vulnerabilidades críticas: redundância de retrabalho entre Google Forms, Autentique, Google Sheets e Notion; campos obrigatórios não preenchidos; permanência indevida de ex-membros como ativos nas listagens de consulta; quebra de fórmulas nas planilhas e risco constante de alteração indevida de pontos disciplinares devido à ausência de restrição de acesso. Além disso, a dependência do WhatsApp para notificar o professor orientador sobre horas e desligamentos despersonaliza o processo institucional e provoca atrasos na formalização acadêmica junto ao SIGAA.

Em resposta a essas dificuldades, a Commit Júnior anseia por uma plataforma web integrada que estabeleça controle unificado em todas as suas frentes de gestão de pessoas. A instituição espera que o ambiente implemente autenticação segura com perfis de Administrador, Jurídico, Líderes de Área e Membros. O cliente deseja a centralização completa do cadastro dos membros com histórico imutável, além do gerenciamento dos fluxos e status de assinatura dos termos de voluntariado, termos de desligamento e certificados corporativos. Almeja-se a automação integral do cômputo de horas de extensão e tempo de casa com base nos cargos exercidos, eliminando as contas manuais e possibilitando o envio formalizado de relatórios diretamente por e-mail para o professor orientador da Commit Júnior realizar os registros no SIGAA. Para a governança disciplinar, busca-se um módulo que automatize a expiração de penalidades e a aplicação de agravantes por reincidência, permitindo também a contestação formal de pontos pendentes pelos membros. Por fim, a organização requer um diretório institucional atualizado que oculte automaticamente membros inativos, mural de avisos por diretoria e calendário corporativo único que viabilize o agendamento de momentos gerais e reuniões individuais (1-on-1) com prevenção de choques de horário.

\clearpage